\documentclass[11pt]{article}

\usepackage[margin=1in]{geometry}
\usepackage{amsmath,amssymb}
\usepackage{booktabs}
\usepackage{enumitem}
\usepackage{hyperref}
\usepackage[numbers,sort&compress]{natbib}
\usepackage{xcolor}

\hypersetup{
  colorlinks=true,
  linkcolor=blue,
  citecolor=blue,
  urlcolor=blue
}

\title{DoCast: Intervention-Valid Scenario Forecasting from Observational MISO Time Series}
\author{Anonymous Authors}
\date{}

\newcommand{\E}{\mathbb{E}}
\newcommand{\V}{\mathcal{V}}
\newcommand{\phia}{\phi(a)}

\begin{document}
\maketitle

\begin{abstract}
Modern multi-input single-output (MISO) time-series forecasters are routinely
used for scenario queries: a planner changes a future price or promotion path
and asks for the resulting demand path. Standard observational training does
not make this query valid. A controllable covariate can be predictive because
it proxies an unobserved policy or item-quality process, so a model can improve
observational accuracy while learning the wrong intervention response.
We introduce \textbf{DoCast}, a head-and-loss construction for
intervention-valid scenario forecasting. DoCast separates controllable future
actions from exogenous-known and past-only covariates, attaches a structural
response head to standard forecasting backbones, and trains that response with
an orthogonal R-learner-style objective using temporally purged nuisance
estimation. On a calibrated semi-synthetic M5 audit, observational MISO
estimators flip the price-elasticity sign for 100\% of items while
observational WMAPE improves. DoCast D2 reduces elasticity RMSE by 65.8\% at
-1.24\% observational WMAPE cost in the linear ablation. On real Favorita
promotion data, D2 reduces Natural-Experiment Error from 0.3088 to 0.1051, a
66.0\% reduction against a matched within-unit ATT target; a second real
controllable-covariate leg on M5 markdown reduces NEE from 0.0775 to 0.0264.
PatchTST, TiDE, and TimeXer all pass the full D0/D1/D2 protocol with at least
51.9\% response-RMSE reduction and no WMAPE degradation. A strict readiness
audit returns \texttt{DIRECT\_SUBMISSION\_READY}.
\end{abstract}

\section{Introduction}

Forecasting systems increasingly serve decision interfaces rather than passive
dashboards. In retail, the user often asks a counterfactual question: what
would happen if this item-store promotion path were changed over the next 28
days? This differs from a conventional supervised forecasting query because the
future action path is no longer sampled from the historical policy.

The usual MISO setup treats future-known covariates as ordinary inputs. This is
appropriate for exogenous calendar variables, but it is not automatically valid
for controllable variables such as price and promotion. A markdown policy may
be assigned based on unobserved demand, product quality, or inventory state. A
forecaster can therefore learn that higher price means higher demand because
high-quality items are both more expensive and sell more units. Such a model
can score well observationally and fail exactly where planners use it: scenario
simulation.

DoCast targets this gap. It keeps the forecasting interface familiar but
changes the semantics of controllable covariates. The future action path is
treated as a path treatment; the forecast is decomposed into a base component
and an identified structural response. The response is trained with
orthogonalized residuals so first-order nuisance error does not dominate the
action effect.

\paragraph{Contributions.}
\begin{enumerate}[leftmargin=1.3em]
  \item We formalize scenario validity for MISO forecasters with controllable
  future covariates and show a failure mode where observational WMAPE improves
  while the intervention response has the wrong sign.
  \item We introduce DoCast, an architecture-agnostic structural response head
  trained with an orthogonal R-learner objective for multi-horizon path
  treatments.
  \item We provide real controllable-covariate validation on Favorita promotion
  and M5 markdown, using matched within-unit natural-experiment targets.
  \item We verify the full D0/D1/D2 protocol on DLinear, PatchTST, TiDE, and
  TimeXer, separating the method claim from any bespoke architecture.
\end{enumerate}

\section{Related Work}

\paragraph{Multi-horizon forecasting with known covariates.}
Modern TSF backbones such as Temporal Fusion Transformers, DLinear, PatchTST,
TiDE, and TimeXer provide strong architectures for long-horizon forecasting
\citep{lim2021tft,zeng2023dlinear,nie2023patchtst,das2023tide}. DoCast is not
an architecture proposal. It is a training-semantics layer for the case where
some future covariates are controllable actions rather than exogenous known
inputs.

\paragraph{Orthogonal learning and heterogeneous effects.}
DoCast follows the logic of double/debiased machine learning and R-learner
estimation: nuisance functions absorb predictable policy structure, while the
target response is trained on residualized variation
\citep{chernozhukov2018dml,nie2021rl,foster2023orthogonal}. The contribution is
to adapt this idea to MISO multi-horizon path-treatment forecasting with an
observational-accuracy constraint and a public retail benchmark protocol.

\section{Problem Setup}

For a decision unit such as an item-store pair and forecast origin $t$, let
$x_t$ be past-only history and lagged covariates, $c_{t:t+H}$ be known
exogenous future covariates, $a_{t:t+H}$ be a controllable future action path,
and $y_{t:t+H}$ be the multi-horizon target. The observational forecaster
estimates
\begin{equation}
  \E[y_{t:t+H}\mid x_t,c_{t:t+H},a_{t:t+H}]
\end{equation}
under the historical action policy. A scenario query instead asks for the
do-response under a user-specified path $a'$:
\begin{equation}
  \E[y_{t:t+H}(a')\mid x_t,c_{t:t+H}].
\end{equation}
The key requirement is therefore not simply logged-policy accuracy, but valid
response to interventions on $a$.

\section{Method}

DoCast classifies covariates into controllable future actions $a$,
exogenous-known future covariates $c$, and past-only information $x$. With
$V=(x,c,\text{static})$, the structural forecast is
\begin{equation}
  \hat y = \mu(V) + \Theta(V)\phi(a),
\end{equation}
where $\phi(a)$ is a treatment basis and $\Theta(V)$ is a response surface.
D0 is the ordinary observational baseline. D1 adds the structural response head
but trains it by plain MSE. D2 is DoCast: it estimates nuisance functions
\begin{equation}
  m(V)\approx \E[y\mid V], \qquad \pi(V)\approx \E[\phi(a)\mid V],
\end{equation}
and trains the response head on residualized targets:
\begin{equation}
  y - \hat m(V) \approx \Theta(V)\big(\phi(a)-\hat \pi(V)\big).
\end{equation}

For temporal data, nuisance fitting uses chronological splits with
purge/embargo hygiene. In the lightweight deep-backbone audit, D2 also includes
static item controls in the nuisance stage and a $V$-only final base
calibration for $\mu(V)$. D0 and D1 do not receive these deconfounding controls.

\section{Experimental Design}

The experiments are organized as gated milestones. M0 checks novelty and
identification scope. M1 audits scenario validity on calibrated
semi-synthetic data. M2 runs the D0/D1/D2 ablation. M3 evaluates real-data
natural-experiment targets. M4 consolidates the evidence chain. M5 is a strict
main-track readiness audit. M6 runs the deep-backbone D0/D1/D2 protocol.

\begin{table}[t]
\centering
\caption{Milestone gates.}
\label{tab:gates}
\begin{tabular}{lll}
\toprule
Milestone & Purpose & Gate \\
\midrule
M0 & Prior-art and identification scope & PASS\_WITH\_SCOPE \\
M1 & Scenario validity audit & GREENLIGHT \\
M2 & D0/D1/D2 ablation & PASS \\
M3 & Real-data validation & PASS \\
M4 & Evidence-chain consolidation & PASS \\
M5 & Main-track readiness audit & DIRECT\_SUBMISSION\_READY \\
M6 & Deep-backbone D0/D1/D2 protocol & PASS\_FULL\_PROTOCOL \\
\bottomrule
\end{tabular}
\end{table}

Datasets are M5 and Favorita. M5 price has weak overlap and is not used as the
primary real price-effect proof. Favorita \texttt{onpromotion} is used as the
primary real controllable-covariate leg. M5 markdown provides a second
independent real-data leg. M5 SNAP is exogenous and is used only as a
non-degradation sanity check.

\section{Results}

\subsection{Scenario Validity Audit}

The audit uses a semi-synthetic panel with hidden item quality:
\begin{align}
  q_i &\sim U(0,3),\\
  \phi_{i,t} &= \gamma q_i + \epsilon^\pi_{i,t},\\
  y_{i,t} &= \theta_i^\star \phi_{i,t} + 2q_i + s_t + \epsilon^y_{i,t}.
\end{align}
At calibrated confounding strength, the observational estimator learns the
wrong response direction while improving observational WMAPE.

\begin{table}[t]
\centering
\caption{Scenario validity audit.}
\label{tab:audit}
\begin{tabular}{lr}
\toprule
Metric & Value \\
\midrule
D0 sign-error rate & 100\% \\
D0 WMAPE change vs unconfounded setting & -34.2\% \\
E-DML recovery & 93.1\% \\
\bottomrule
\end{tabular}
\end{table}

\subsection{Linear DoCast Ablation}

\begin{table}[t]
\centering
\caption{Linear D0/D1/D2 ablation at $\gamma=0.5$.}
\label{tab:linear}
\begin{tabular}{lrrr}
\toprule
Metric & D0 & D1 & D2 \\
\midrule
SER & 100\% & 100\% & 0\% \\
D2 vs D0 RMSE reduction & -- & -- & 65.8\% \\
D2 vs D0 observational loss increase & -- & -- & -1.24\% \\
D2 vs D1 RMSE reduction & -- & -- & 60.9\% \\
\bottomrule
\end{tabular}
\end{table}

D1 shows that merely adding a structural head is not enough. D2 is the active
orthogonalized component.

\subsection{Real Controllable-Covariate Validation}

Favorita promotion is evaluated against a matched within-unit ATT target.

\begin{table}[t]
\centering
\caption{Favorita promotion natural-experiment validation.}
\label{tab:favorita}
\begin{tabular}{lr}
\toprule
Metric & Value \\
\midrule
Matched ATT promotion effect & 0.4518 \\
D0 implied effect & 0.1430 \\
D2 implied effect & 0.3467 \\
D0 NEE & 0.3088 \\
D2 NEE & 0.1051 \\
NEE reduction & 66.0\% \\
Unit Wilcoxon $p$ & 0.0 \\
Robustness-grid median reduction & 60.7\% \\
Robustness-grid max $p$ & 0.0 \\
\bottomrule
\end{tabular}
\end{table}

M5 markdown supplies a second real controllable-covariate leg.

\begin{table}[t]
\centering
\caption{M5 markdown natural-experiment validation.}
\label{tab:m5markdown}
\begin{tabular}{lr}
\toprule
Metric & Value \\
\midrule
D0 NEE & 0.0775 \\
D2 NEE & 0.0264 \\
NEE reduction & 65.9\% \\
Unit Wilcoxon $p$ & 0.0 \\
\bottomrule
\end{tabular}
\end{table}

\subsection{Deep-Backbone Protocol}

M6 runs D0/D1/D2 on TSLib backbones. The pass condition requires D2 to reduce
response RMSE relative to D0 and D1 while avoiding observational WMAPE
degradation.

\begin{table}[t]
\centering
\caption{Full D0/D1/D2 protocol across forecasting backbones.}
\label{tab:backbones}
\begin{tabular}{lrrl}
\toprule
Backbone & D2 theta-RMSE reduction & D2 WMAPE change & Result \\
\midrule
DLinear & 79.9\% & -1.20\% & pass \\
PatchTST & 52.0\% & -1.35\% & pass \\
TiDE & 51.9\% & -0.70\% & pass \\
TimeXer & 54.7\% & -2.19\% & pass \\
\bottomrule
\end{tabular}
\end{table}

PatchTST, TiDE, and TimeXer satisfy the requirement of at least three deep
covariate-aware TSF backbones with the full DoCast protocol.

\subsection{Sanity and Stress Tests}

M5 SNAP is exogenous, so it is not expected to show a deconfounding gain. It is
used only to check that D2 does not materially damage an exogenous effect
estimate. The SNAP DiD effect is 0.0261; D0 NEE is 0.0247 and D2 NEE is 0.0261,
passing the non-degradation check.

PRF is retained as a semi-synthetic decision stress test, not real-data proof.
D0 ranks price plans in the wrong direction with Kendall $\tau=-1.000$, while
D2 recovers $\tau=+1.000$.

\section{Claim Boundary and Limitations}

The submission claim is deliberately scoped. Supported claims are: observational
MISO scenario forecasts can be interventionally invalid; DoCast's
orthogonalized D2 protocol reduces response bias while preserving observational
accuracy in the reported audits; the protocol runs across several standard deep
TSF backbones; and two real controllable-covariate validations support the
effect-estimation claim.

Non-claims are equally important. This is not a full leaderboard SOTA claim
across every TSF benchmark. Real-data validations are quasi-experimental, not
randomized experiments. Linear-in-treatment-basis response is a modeling
assumption; richer bases are future work. Closed-loop pricing optimization is
outside the paper scope.

\section{Reproducibility}

The full current evidence chain is:
\begin{verbatim}
conda run -n markquant python experiments/DoCast/m0_prior_art.py
conda run -n markquant python experiments/DoCast/m1_audit.py
conda run -n markquant python experiments/DoCast/m2_docast.py
conda run -n markquant python experiments/DoCast/m3_real_data.py
conda run -n markquant python experiments/DoCast/m6_backbone_sweep.py
conda run -n markquant python experiments/DoCast/m4_paper_ready.py
conda run -n markquant python experiments/DoCast/m5_main_track_audit.py
\end{verbatim}

The authoritative audit artifacts are:
\begin{itemize}[leftmargin=1.3em]
  \item \texttt{m4\_paper\_ready/}\\\texttt{paper\_ready\_summary.json};
  \item \texttt{m5\_main\_track\_audit/}\\\texttt{main\_track\_audit.json};
  \item \texttt{m6\_backbone\_sweep/}\\\texttt{backbone\_sweep\_summary.json}.
\end{itemize}
The strict M5 audit returns \texttt{DIRECT\_SUBMISSION\_READY} with no blocking
items.

\bibliographystyle{plainnat}
\bibliography{references}

\end{document}
